% Generated by scripts/assemble.py; corrections belong in source/corrections.json.
\examhead{38}{华东师范大学 2026 年数学分析试卷}{41}
% p042-b12; source PDF page 42
\par\Needspace{5\baselineskip}
\blockanchor{p042-b12}
\majorquestion{一}填空题(每题8分, 共56分)\par
% p042-b13; source PDF page 42
\blockanchor{p042-b13}
\question{1}计算 $\sum_{n=1}^{\infty} \frac{2n-1}{2^n} =$ \_\_\_\_.\par
% p042-b14; source PDF page 42
\blockanchor{p042-b14}
\question{2}计算  $\lim_{x\to0^{+}}\sqrt[x]{\sec\sqrt{2x}}=$  \_\_\_\_.\par
% p042-b15; source PDF page 42
\blockanchor{p042-b15}
\question{3}计算 $\int_0^1\int_0^{\sqrt{1 - x^2}}\int_{\sqrt{x^2 + y^2}}^{\sqrt{2 - x^2 - y^2}}xy\mathrm{d}z\mathrm{d}y\mathrm{d}x =$\par
% p042-b16; source PDF page 42
\blockanchor{p042-b16}
\question{4}设 $|a| \neq 1$, 计算 $\int_{0}^{\frac{\pi}{2}} \frac{\arctan(a \tan x)}{\tan x} \, \mathrm{d}x =$ \_\_\_\_.\par
% p042-b17; source PDF page 42
\blockanchor{p042-b17}
\question{5}求和函数 $\sum_{n=1}^{\infty} \frac{(-1)^{n-1}}{n(n+1)} (x^2 - 1)^{2n} =$ \_\_\_\_.\par
% p042-b18; source PDF page 42
\blockanchor{p042-b18}
\question{6}计算第二型曲面积分 $\iint_{\Sigma} x z^{2} \, \mathrm{d}y \, \mathrm{d}z + \frac{1}{y} \, \mathrm{d}z \, \mathrm{d}x + x^{2} z \, \mathrm{d}x \, \mathrm{d}y =$ \_\_\_\_, 其中 $\Sigma$ 为抛物面 $y = x^{2} + z^{2}$ 被平面 $y = 1, y = 2$ 所截部分, 取外侧.\par
% p042-b19; source PDF page 42
\blockanchor{p042-b19}
\question{7}设 $0 < x < 2\pi, a \neq 0$, 则 $\frac{1}{2a} + \sum_{n=1}^{\infty} \frac{a \cos nx - n \sin nx}{n^2 + a^2} =$ \_\_\_\_.\par
% p042-b20; source PDF page 42
\par\Needspace{5\baselineskip}
\blockanchor{p042-b20}
\majorquestion{二}解答题(需提供以下各题的详细解答过程)\par
% p042-b21; source PDF page 42
\blockanchor{p042-b21}
\question{8}(15 分) 解答如下问题:\par
% p042-b22; source PDF page 42
\blockanchor{p042-b22}
(1) 证明: 对任意正整数 $n$, 方程 $x^{3} + 3x + \frac{1}{n} = 0$ 有且仅有一个实数解.\par
% p042-b23; source PDF page 42
\blockanchor{p042-b23}
(2) 设 $x_{n}, n = 1,2,\dots$ 为 (1) 中方程的实根, 证明: 数列 $\{x_{n}\}$ 收敛.\par
% p042-b24; source PDF page 42
\blockanchor{p042-b24}
\question{9}(15 分) 设 $\{f_n(x)\}$ 为定义在 $[a, b]$ 上的函数列, 对任意 $x \in [a, b]$, 数列 $\{f_n(x)\}$ 有界, 且存在 $K > 0$ 使得对任意 $x, y \in [a, b]$ 和 $n \in \mathbb{N}_+$, 都有 $|f_n(x) - f_n(y)| \leq K|x - y|$. 证明: 函数列 $\{f_n(x)\}$ 必存在一致收敛的子列.\par
% p042-b25; source PDF page 42
\blockanchor{p042-b25}
\question{10}(16 分) 判断下列命题是否正确, 若正确给出证明, 若错误给出反例说明.\par
% p043-b00; source PDF page 43
\blockanchor{p043-b00}
(1) 若 $f(x)$ 是 $[a, b]$ 上的凸函数, 则 $f(x)$ 在 $[a, b]$ 上连续.\par
% p043-b01; source PDF page 43
\blockanchor{p043-b01}
(2) 若 $f(x)$ 是 $(a, b)$ 上的凸函数, 且 $f(x)$ 在 $(a, b)$ 上可微, 则 $f'(x)$ 在 $(a, b)$ 上连续.\par
% p043-b02; source PDF page 43
\par\Needspace{6\baselineskip}
\blockanchor{p043-b02}
\question{11}(16 分) 设 $f(x)$ 在 $\left[0, \frac{\pi}{2}\right]$ 上为正值连续函数, 且\par
\nopagebreak[4]
% p043-b03; source PDF page 43
\blockanchor{p043-b03}
\displayformula{I _ {1} = \int_ {0} ^ {\frac {\pi}{2}} e ^ {f (x)} \sin x \ln f (x) \mathrm{d} x = 0, I _ {2} = \int_ {0} ^ {\frac {\pi}{2}} e ^ {f (x)} \cos x \ln f (x) \mathrm{d} x = 0.}
% p043-b04; source PDF page 43
\blockanchor{p043-b04}
证明：\par
% p043-b05; source PDF page 43
\blockanchor{p043-b05}
(1) 在 $\left(0, \frac{\pi}{2}\right)$ 上至少存在 $x_{1}$, 使得 $f(x_{1}) = 1$.\par
% p043-b06; source PDF page 43
\blockanchor{p043-b06}
(2) 在 $\left(0, \frac{\pi}{2}\right)$ 上至少存在 $x_{2}, x_{3}$, 使得 $f(x_{2}) = f(x_{3})$.\par
% p043-b07; source PDF page 43
\blockanchor{p043-b07}
\question{12}(16 分) 设二元函数 $f(x, y)$ 的偏导数 $f_x, f_y, f_{xy}$ 在点 $P_0(x_0, y_0)$ 处都连续. 证明: $f_{yx}(x_0, y_0)$ 存在, 且 $f_{yx}(x_0, y_0) = f_{xy}(x_0, y_0)$.\par
% p043-b08; source PDF page 43
\blockanchor{p043-b08}
\question{13}(16 分) 设 $\Omega$ 是 $\mathbb{R}^3$ 中的有界单连通区域, 其边界 $S$ 是一个光滑封闭曲面, $\vec{n}$ 是 $S$ 上动点 $Q(\xi, \eta, \zeta)$ 处的单位外法向量, $P_0(x_0, y_0, z_0)$ 为不在 $S$ 上的一点, $r = \sqrt{(\xi - x_0)^2 + (\eta - y_0)^2 + (\zeta - z_0)^2}$, 证明: \displayformula{\iiint_{\Omega} \frac{1}{r} \mathrm{d}\xi \mathrm{d}\eta \mathrm{d}\zeta = \frac{1}{2} \iint_{S} \cos \angle (\overrightarrow{PQ}, \overrightarrow{n}) \mathrm{d}S}, 其中 $\angle (\overrightarrow{PQ}, \overrightarrow{n})$ 表示向量 $\overrightarrow{PQ}$ 和 $\overrightarrow{n}$ 的夹角.\par
